\subsection{Objekt T5}
\label{subsec:ObjektT5}


\small

\paragraph{1. Objektsteckbrief}

\par Bei dem Objekt T5 handelt es sich um ein Wohn- und Geschäftshaus in der Tannenstraße 5, 92334 Pollanten. Es handelt sich um ein im Jahr 2003 errichtetes Gebäude, das laufend modernisiert und instand gehalten wurde. Das Objekt verfügt über 5 Wohneinheiten und 1 Gewerbeeinheit mit einer vermietbaren Fläche von insgesamt [\(m^2\) Wohnen] + [\(m^2\) Gewerbe]. Die Vermietungsquote liegt konstant bei 100\%. Das Objekt befindet sich im Eigentum der \gls{stz} und stellt eine wichtige Komponente des bestehenden Immobilienportfolios dar.


\begin{table} [htbp!]
\begin{tabular*}{\textwidth}{@{\extracolsep{\fill}}>{\raggedright\arraybackslash}p{4.5cm}>{\raggedright\arraybackslash}p{10.3cm}}
\toprule
\textbf{Merkmal} & \textbf{Angabe} \\
\midrule
Adresse & Tannenstraße 5, 92334 Pollanten \\
Objektart & Wohn- und Geschäftshaus \\
Baujahr / Modernisierung & 2003 / laufend \\
Vermietbare Fläche & [\(m^2\) Wohnen] + [\(m^2\) Gewerbe] \\
Einheiten & 5, 1 \\
Vermietungsstand & 100\% per \today \\
Eigentümer & \gls{stz} \\
\bottomrule
\end{tabular*}
\caption{Objektsteckbrief T5}
\label{tab:objektsteckbriefT5}
\end{table}

\paragraph{2. Ertragslage und Stabilität (p.a.)}
\begin{tabular*}{\textwidth}{@{\extracolsep{\fill}}>{\raggedright\arraybackslash}p{8.4cm}>{\raggedleft\arraybackslash}p{3.1cm}>{\raggedleft\arraybackslash}p{3.1cm}}
\toprule
\textbf{Kennzahl} & \textbf{Ist} & \textbf{Stabilisiert} \\
\midrule
Nettokaltmiete & [EUR] & [EUR] \\
Nicht umlagefähige Bewirtschaftungskosten & [EUR] & [EUR] \\
NOI (Net Operating Income) & [EUR] & [EUR] \\
Leerstandsquote & 0\% & 15\% \\
\bottomrule
\end{tabular*}

\paragraph{3. Bankrelevante Kennzahlen (stabilisierter Zustand)}
\begin{tabular*}{\textwidth}{@{\extracolsep{\fill}}>{\raggedright\arraybackslash}p{6.1cm}>{\raggedleft\arraybackslash}p{2.8cm}>{\raggedright\arraybackslash}p{6.3cm}}
\toprule
\textbf{Kennzahl} & \textbf{Wert} & \textbf{Einordnung} \\
\midrule
Marktwert / interner Ansatz & [EUR] & Bewertungsstichtag: [Datum] \\
Beleihungswert (konservativ) & [EUR] & Sicherheitsabschlag: [\%] \\
Wunschdarlehen auf Objektbasis & [EUR] & Anteil am Gesamtobligo: [\%] \\
LTV (Darlehen / Marktwert) & [\%] & Bankziel: [z.B. \textless= 70\%] \\
Debt Yield (NOI / Darlehen) & [\%] & Bankziel: [z.B. \textgreater= 8\%] \\
DSCR (CFADS / Schuldendienst) & [x,xx] & Bankziel: [z.B. \textgreater= 1,20] \\
\bottomrule
\end{tabular*}

\paragraph{4. Herleitung des konservativen Kreditrahmens}
\begin{tabular*}{\textwidth}{@{\extracolsep{\fill}}>{\raggedright\arraybackslash}p{10.8cm}>{\raggedleft\arraybackslash}p{3.4cm}}
\toprule
\textbf{Limitierende Sicht} & \textbf{Darlehenslimit} \\
\midrule
LTV-basierter Rahmen & [EUR] \\
DSCR-basierter Rahmen & [EUR] \\
Debt-Yield-basierter Rahmen & [EUR] \\
\midrule
\textbf{Empfohlener Kreditrahmen (Minimum aus den drei Sichten)} & \textbf{[EUR]} \\
\bottomrule
\end{tabular*}

\paragraph{5. Risiko- und Sicherheitenprofil}
\begin{itemize}
	\item \textbf{Mietausfall-/Leerstandsrisiko:} [kurze Bewertung], abgesichert durch [Maßnahme, z.B. Mietreserve in EUR].
	\item \textbf{Instandhaltungsrisiko:} [kurze Bewertung], abgesichert durch [Capex-Plan / Rücklage].
	\item \textbf{Marktrisiko:} [kurze Bewertung], mitigiert durch [Lagequalität / diversifizierte Mieterstruktur].
	\item \textbf{Rechtliche Sicherheiten:} erstrangige Grundschuld [EUR], Abtretung Mietforderungen, weitere Covenants gemäß Kreditvertrag.
\end{itemize}

\paragraph{6. Kurzfazit für die Kreditentscheidung}
Das Bestandsobjekt Tannenstraße 5 zeigt im stabilisierten Zustand eine tragfähige Cashflow-Basis zur Bedienung von Zins und Tilgung. Unter konservativen Annahmen ergibt sich ein belastbarer Kreditrahmen von \textbf{[EUR]}, der durch werthaltige Sicherheiten und ein transparentes Risikomanagement unterlegt ist.

\vspace{0.5em}
\textit{Hinweis: Detailrechnungen (Tilgungsplan, Sensitivitäten, Mieterliste) sind im Anhang bzw. in den Tabellenblättern hinterlegt.}

\normalsize
